\documentclass[aspectratio=169]{beamer}
\usepackage{tikz}
\usepackage{graphicx}
\usepackage{listings}
\usepackage{xcolor}
\usepackage{hyperref}
\usepackage[spanish,es-noquoting]{babel}
\usepackage[round]{natbib}

\lstset{
  basicstyle=\ttfamily\small,
  keywordstyle=\color{blue},
  stringstyle=\color{red},
  commentstyle=\color{gray},
  showstringspaces=false
}
% Minimal look
\usetheme{Warsaw}
%\setbeamertemplate{headline}
\setbeamertemplate{footline}{
  \leavevmode%
  \hbox{%
    % Left: Section name
    \begin{beamercolorbox}[wd=.33\paperwidth,ht=2.25ex,dp=1ex,left]{author in head/foot}%
      \hspace*{2ex}\insertsection
    \end{beamercolorbox}%
    % Center: Presentation title
    \begin{beamercolorbox}[wd=.34\paperwidth,ht=2.25ex,dp=1ex,center]{title in head/foot}%
      \insertshorttitle
    \end{beamercolorbox}%
    % Right: Slide numbers
    \begin{beamercolorbox}[wd=.33\paperwidth,ht=2.25ex,dp=1ex,right]{date in head/foot}%
      \insertframenumber{} / \inserttotalframenumber\hspace*{2ex}
    \end{beamercolorbox}%
  }%
  \vskip0pt%
}

\usecolortheme{default}
\usefonttheme{serif}
\graphicspath{{../}{./}}
\setbeamertemplate{caption}[numbered]

% Title info
\title[Cooperación Internacional y Pr]{Cooperación Internacional y Producción Artística en Colombia}
\subtitle{Una Perspectiva desde la Intermediación Cultural Experta}
\author{Maestro Diseñador Instruccional (Experto Intermediario Cultural)}
\institute{Su Institución / Nuestro Equipo}
\date{\today}



\begin{document}

% Title slide
\begin{frame}
  \titlepage
\end{frame}

% Reset background for normal slides
\usebackgroundtemplate{}



\begin{frame}[fragile]{Cooperación Internacional y Producción Artística en Colombia}


\end{frame}



\begin{frame}[fragile]{1. Objetivos de la Sesión y Contexto General}


\textbf{Objetivos de la Sesión}

\begin{itemize}
  \item \textbf{Comprender} la influencia de la cooperación internacional en la producción artística colombiana.
  \item \textbf{Analizar} los marcos conceptuales clave y las dinámicas inherentes a esta relación.
  \item \textbf{Identificar} los desafíos y oportunidades desde la perspectiva de un intermediario cultural experto.
  \item \textbf{Fomentar} la reflexión crítica sobre la autonomía y sostenibilidad artística en contextos de cooperación.
\end{itemize}
\textbf{Contexto General}

Colombia, un país de vasta riqueza cultural y diversidad artística, ha sido históricamente un receptor y, en creciente medida, un actor en la cooperación internacional. Esta relación ha impactado significativamente su producción artística, trascendiendo lo puramente financiero para incidir en narrativas, estéticas y procesos creativos.

\end{frame}



\begin{frame}[fragile]{2. Marco Conceptual y Conceptos Clave I}


\textbf{Cooperación Internacional y Producción de Significados}

\begin{itemize}
  \item \textbf{Más allá del soporte financiero:} La cooperación internacional no es solo una transferencia de recursos económicos.
  \item \textbf{Moldeando narrativas culturales:} Influye profundamente en la ``producción de significados'' (Cajiao Restrepo, 1999) dentro de las obras de arte.
  \item \textbf{Intercambio cultural profundo:} Implica la transmisión de valores estéticos, metodologías artísticas y marcos narrativos, impactando la interpretación de la realidad a través del arte.
  \item \textbf{Impacto en la identidad y representación:} Afecta cómo se construyen y presentan las identidades y realidades culturales colombianas en la esfera artística.
\end{itemize}
\end{frame}



\begin{frame}[fragile]{3. Marco Conceptual y Conceptos Clave II}


\textbf{La ``Cooperación'' bajo el Microscopio}

\begin{itemize}
  \item \textbf{Análisis crítico:} Es fundamental examinar la verdadera naturaleza de la ``cooperación'' en los proyectos culturales (Herrera, 2017).
  \item \textbf{Dinámicas de poder:} Cuestiona si las iniciativas internacionales promueven una autonomía artística genuina o si, inadvertidamente, generan dependencias.
  \item \textbf{Sostenibilidad y agendas:} Se pregunta si los proyectos se alinean con las necesidades artísticas locales o si responden a prioridades dictadas por los donantes.
  \item \textbf{Rol del intermediario cultural:} Permite escudriñar la equidad de las asociaciones y asegurar que la producción cultural sea mutuamente beneficiosa.
\end{itemize}
\end{frame}



\begin{frame}[fragile]{4. Marco Conceptual y Conceptos Clave III}


\textbf{Colombia en el Escenario Internacional}

\begin{itemize}
  \item \textbf{Contexto de política exterior:} La participación de Colombia con la comunidad internacional enmarca las iniciativas de cooperación cultural (Tassara, 2015).
  \item \textbf{Actores internacionales:} Identifica a los principales agentes (gobiernos, ONGs, fundaciones) y sus intereses estratégicos.
  \item \textbf{Marcos de desarrollo:} La financiación y las oportunidades para el sector artístico a menudo están ligadas a objetivos diplomáticos, esfuerzos de construcción de paz o estrategias de desarrollo económico más amplias.
  \item \textbf{Influencia en prioridades temáticas:} Estos marcos impactan qué tipos de arte se apoyan y qué narrativas se promueven.
\end{itemize}
\end{frame}



\begin{frame}[fragile]{5. Casos de Estudio y Ejemplos Prácticos Aplicados I}


\textbf{Influencia en Narrativas y Estéticas Artísticas}

\begin{itemize}
  \item \textbf{Hibridación cultural:} La cooperación internacional frecuentemente conduce a expresiones artísticas híbridas, fusionando elementos locales con influencias estéticas globales.
  \item \textbf{Transferencia de metodologías:} Se observa la adopción de nuevas técnicas, enfoques curatoriales y formatos de exhibición procedentes de otros contextos culturales.
  \item \textbf{Modernización y estandarización:} Puede implicar una ``modernización'' del sector artístico, alineando prácticas profesionales y estándares organizacionales con modelos internacionales.
  \item \textbf{Desafío a la distinción local:} Esto puede generar el reto de preservar la singularidad y la pertinencia local frente a una sutil homogeneización cultural.
\end{itemize}
\end{frame}



\begin{frame}[fragile]{6. Casos de Estudio y Ejemplos Prácticos Aplicados II}


\textbf{Desafíos en la Práctica de la ``Cooperación Genuina''}

\begin{itemize}
  \item \textbf{Equidad en las alianzas:} Asegurar que los proyectos culturales sean colaborativos de manera auténtica y no una mera importación de modelos extranjeros.
  \item \textbf{Alineación de prioridades:} El desafío de conciliar las agendas de los donantes con las necesidades y visiones de los artistas y gestores culturales colombianos.
  \item \textbf{Fomento de la autonomía artística:} Evitar la creación de dependencias y empoderar a las instituciones y creadores locales para que lideren sus propios procesos.
  \item \textbf{Navegación de corrientes geopolíticas:} El intermediario cultural debe equilibrar la obtención de recursos externos con la defensa de la libertad creativa y la integridad cultural local.
\end{itemize}
\end{frame}



\begin{frame}[fragile]{7. Ejercicios Prácticos y Actividades Interactivas I}


\textbf{Análisis de un Escenario Hipotético}

\textbf{Actividad: ``El Proyecto del Taller Urbano''}

\textbf{Escenario:}
Una fundación internacional financia un proyecto de arte urbano en un barrio vulnerable de Bogotá. El objetivo declarado es ``promover la expresión juvenil y la cohesión social''. La fundación insiste en utilizar un estilo muralista figurativo popular en Europa y un formato de taller que prioriza la técnica sobre la narrativa local. Los artistas locales, aunque agradecidos por los recursos, sienten que sus propias estéticas y temáticas son relegadas.

\textbf{Pregunta para la audiencia (discusión grupal de 5 minutos):}

\begin{itemize}
  \item ¿Qué elementos de `\texttt{producción de significados'' o }`currículo invisible'' identifican en este escenario?
  \item ¿Dónde residen las dinámicas de poder en esta ``cooperación''?
\end{itemize}
\end{frame}



\begin{frame}[fragile]{8. Ejercicios Prácticos y Actividades Interactivas II}


\textbf{Rol del Intermediario Cultural}

\textbf{Actividad: ``Estrategias para la Sostenibilidad Cultural''}

\textbf{Instrucciones (trabajo en parejas de 10 minutos):}

\begin{itemize}
  \item \textbf{Imaginen} que son intermediarios culturales expertos encargados de renegociar el proyecto ``El Taller Urbano''.
  \item \textbf{Diseñen} dos estrategias concretas para transformar esta iniciativa en una cooperación más equitativa y autónoma para los artistas locales.
  \begin{itemize}
    \item Consideren cómo balancear la financiación internacional con la promoción de la identidad artística local.
  \end{itemize}
  \item \textbf{Compartan} brevemente sus propuestas con el grupo.
\end{itemize}
\end{frame}



\begin{frame}[fragile]{9. ``Sabías que...'' / Datos Curiosos}


\textbf{Perspectivas Sorprendentes de la Investigación}

\begin{itemize}
  \item \textbf{El ``currículo invisible'' de la ayuda cultural:} Más allá de los objetivos explícitos, la cooperación internacional a menudo introduce prácticas organizativas, estándares profesionales e incluso preferencias estéticas de los países donantes, influyendo en la formación y exhibición artística en Colombia.
  \item \textbf{Un nexo poco investigado:} Existe una relativa escasez de literatura académica directa y sistemática que conecte explícitamente `\texttt{cooperación internacional'' con }`producción artística'' en Colombia. Esto sugiere una brecha significativa en la investigación sobre una dinámica crucial para el arte contemporáneo colombiano.
  \item \textbf{Dinámicas de poder y agencia:} A pesar de la fachada benévola de la `\texttt{cooperación'', las asimetrías de poder son inherentes. La pregunta }`¿Cuánto de cooperación hay en la 'cooperación internacional'?'' (Herrera, 2017) es central, subrayando el desafío de fomentar la verdadera agencia y autodeterminación en la producción artística colombiana.
\end{itemize}
\end{frame}



\begin{frame}[fragile]{10. Conclusiones y Puntos Clave}


\textbf{Takeaways de la Sesión}

\begin{itemize}
  \item \textbf{Catalizador complejo:} La cooperación internacional es un motor multifacético que impulsa y, a la vez, desafía la producción artística colombiana, incidiendo en sus narrativas, estéticas y estructuras.
  \item \textbf{Enfoque crítico proactivo:} Es imperativo acercarse a la cooperación con una lente crítica, comprendiendo sus dinámicas de poder y sus implicaciones más allá de lo financiero.
  \item \textbf{Rol vital del intermediario cultural:} La figura del intermediario cultural experto es esencial para negociar, contextualizar y asegurar que las alianzas sean equitativas, sostenibles y respetuosas de la autonomía artística local.
  \item \textbf{Potencial para la innovación y sostenibilidad:} Cuando se gestiona estratégicamente, la cooperación internacional puede ser una fuente invaluable de recursos, conocimientos y oportunidades para el florecimiento y la proyección global del arte colombiano, fortaleciendo la identidad y la diversidad cultural.
\end{itemize}
\end{frame}



\begin{frame}[fragile]{11. Referencias Bibliográficas}


\begin{itemize}
  \item Cajiao Restrepo, F. (1999). Cooperación internacional y producción de significados: todo lo que lleva dentro una hamburguesa. \textit{Colombia Internacional}, (46), 25-34. DOI: 10.7440/colombiaint46.1999.02
  \item Herrera, J. (2017). ¿Cuánto de cooperación hay en la “cooperación internacional”?. In \textit{Cooperación y academia: una relación pendiente en Colombia} (pp. 165-178). Cali, Colombia: Universidad Icesi. DOI: 10.2307/j.ctv18msn54.12
  \item Tassara, C. (2015). Colombia en el escenario internacional. \textit{Revista Internacional de Cooperación y Desarrollo}, 2(1), 1-13. DOI: 10.21500/23825014.2279
\end{itemize}
\end{frame}




\begin{frame}[allowframebreaks]{References}
  \bibliographystyle{apalike}
  \bibliography{references}
\end{frame}


\end{document}